일반적으로 PC는 다목적 시스템인데 비해 임베디드 시스템은 흔히 내장형 시스템이라고 하며 정해진 용도에 맞추어서 최적화되어 있다. 임베디드 시스템은 PC에서와는 다른 개발 플랫폼을 사용하기 때문에 많은 애로사항이 발생한다. 현재의 임베디드 환경에 사용되어지는 마이크로프로세서의 종류도 다양하며 소프트웨어 플랫폼 역시 무수히 많은 임베디드 OS가 사용된다.

본 논문에서는 기존에 개발되어서 사용되고 있는 임베디드 OS 중에서 공개 소프트웨어인 임베디드 리눅스 시스템을 이용하여 향상된 기능의 라우터를 구현한다. 임베디드 리눅스를 사용하는 장점으로는 네트워크 호환성이 뛰어나며, 크기가 작으며, 쉽게 재구성화가 가능하다는 점이다. 본 논문에서는 구현과정에 대한 자세한 기술보다는 전체적인 작업과 사용자 인터페이스 부분에 초점을 맞추어서 기술한다.